\PuzzleChapterIntro{The Five Seals of the Linear Vault}{Determinants \textbullet\ Interlocked Gates \textbullet\ The Vault Exponent}{This is a chain of determinant problems where parameters cascade from one matrix to the next. You must compute the 2-adic valuation (the power of 2) of the final determinant.}

\Lore{
At the last door of the Linear Vault, Eryn finds a mercy carved into the lintel:

\begin{quote}
\textit{Do not speak the whole verdict. Speak only its height.}
\end{quote}

Five seals must still be set—staircase, mirror, choir, black--white wall, and the imaginary bolt—
but the tribute is not the determinant itself.

The Vault asks only for how many times the verdict can be halved before an odd number remains.
}

\Problem
Athena must forge the \emph{Spine Matrix} $\Omega$ from five interlocked seals.

\noindent\textbf{(I) The Staircase Seal.}
Let
\[
J=\begin{pmatrix}
1&1&0&0\\
0&1&1&0\\
0&0&1&1\\
0&0&0&1
\end{pmatrix}.
\]
Let \(u\) be the \((1,2)\) entry of \(J^5\).

\noindent\textbf{(II) The Mirror Seal.}
Using this \(u\), define the \(4\times 4\) matrix
\[
L(u)=\begin{pmatrix}
2u&2&u&u\\
u&1&1&1\\
2&2&u^2&u\\
1&1&u&1
\end{pmatrix}.
\]

\noindent\textbf{(III) The Choir Seal.}
Set \(n=u-2\) and \(\alpha=\dfrac{\pi}{n+1}\). Define the \(n\times n\) matrix
\[
A_{m,j}=\sin(mj\alpha)\qquad(1\le m,j\le n),
\]
and the \(2n\times 2n\) matrix
\[
S=\begin{pmatrix}
0&A\\
A&0
\end{pmatrix}.
\]

\noindent\textbf{(IV) The Black--White Wall.}
Let \(m=2u+1\) and define the \(m\times m\) matrix \(W\) by
\[
W_{i,j}=\begin{cases}
1,& i\le j,\\
-1,& i>j.
\end{cases}
\]

\noindent\textbf{(V) The Imaginary Bolt.}
Define the \(2m\times 2m\) matrix
\[
C=\begin{pmatrix}
W&-W\\
W&W
\end{pmatrix}.
\]

\noindent\textbf{The Spine.}
Let
\[
\Omega=\mathrm{diag}\bigl(L(u),\,S,\,C\bigr).
\]

\noindent\textbf{Vault Format.}
There is a unique integer \(E\ge 0\) and a unique odd integer \(Q\) such that
\[
\det(\Omega)=Q\cdot 2^{E}.
\]

\VaultDirective{Output $E$, the height of the verdict.}

\SolutionPage
\(\Omega\) is block diagonal, so
\[
\det(\Omega)=\det(L(u))\det(S)\det(C),
\]
and \(E=\nu_2(\det L)+\nu_2(\det S)+\nu_2(\det C)\).

\smallskip
\textbf{Staircase.} \(J=I+N\) with \(N^4=0\). Then \(J^5=\sum_{k=0}^3\binom{5}{k}N^k\), so the \((1,2)\) entry is \(\binom51=5\). Thus \(u=5\), hence \(n=u-2=3\), \(m=2u+1=11\).

\smallskip
\textbf{Mirror.} The built-in identity \(\det(L(u))=(u-1)^2(u-2)^2\) gives
\[
\det(L(5))=4^2\cdot 3^2=9\cdot 2^4\ \Rightarrow\ \nu_2=4.
\]

\smallskip
\textbf{Choir.} For \(n=3\), \(\alpha=\pi/4\). Standard sine orthogonality yields \(A^2=\frac{n+1}{2}I=2I\), so \(\det(A)^2=\det(2I_3)=2^3\).
Also \(\det\!\begin{psmallmatrix}0&A\\A&0\end{psmallmatrix}=(-1)^n\det(A)^2\), hence \(\det(S)=-2^3\) and \(\nu_2=3\).

\smallskip
\textbf{Bolt.} For the “wall” matrix \(W\) of size \(m=11\), one has \(\det(W)=2^{m-1}=2^{10}\) (row-add then triangularize).
Now
\[
C=\begin{pmatrix}W&-W\\W&W\end{pmatrix}\ \xrightarrow{\text{row }2\leftarrow 2+1}\
\begin{pmatrix}W&-W\\2W&0\end{pmatrix},
\]
so factoring \(2\) from the bottom \(m\) rows gives \(\det(C)=2^{m}\det\!\begin{psmallmatrix}W&-W\\W&0\end{psmallmatrix}\).
Since \(W\) is invertible, the Schur complement gives \(\det\!\begin{psmallmatrix}W&-W\\W&0\end{psmallmatrix}=\det(W)^2\).
Thus
\[
\det(C)=2^{11}\cdot (2^{10})^2=2^{31}\ \Rightarrow\ \nu_2=31.
\]

Therefore \(E=4+3+31=38\).

\FinalAnswer{38}

\NextPage
